%!TeX root=MemoriaTFG.tex

\chapter{Documento de entrega del dataset DermapixelAI 1.0}
\label{anexo:dataset-entrega}
\label{cap:anexoe}

Este anexo recoge el documento de entrega formal del dataset
DermapixelAI 1.0, complementario al anexo~\ref{anexo:pipeline} que
documenta su pipeline interno de construcción. El presente anexo
sigue el formato \emph{Datasheet for Datasets} propuesto
por~\cite{datasheets2018} y reúne la información necesaria para la
reutilización académica del dataset.

\section{Identidad del dataset}
\label{sec:dataset-identidad}

\begin{itemize}
\item \textbf{Nombre:} DermapixelAI.
\item \textbf{Versión:} 1.0 (primera versión estable y publicable).
\item \textbf{Fecha de cierre de la versión 1.0:} mayo de 2026.
\item \textbf{Identificador persistente (DOI):} pendiente de
asignación en el momento de la publicación efectiva del dataset.
El depósito formal se realiza en \textbf{Zenodo}
(\url{https://zenodo.org}) bajo licencia CC~BY-NC-SA~4.0 con DOI
permanente y citable, complementado con el repositorio de la
Universitat de les Illes Balears cuando éste habilite el alojamiento
de datasets binarios de tamaño equivalente.
\item \textbf{Cita corta:} \emph{Taberner, R. y Contestí Coll, A.
(2026). DermapixelAI 1.0: dataset dermatológico clínico en
castellano. Versión 1.0.}
\end{itemize}

\section{Procedencia y autoría}
\label{sec:dataset-procedencia}

El dataset DermapixelAI 1.0 se construye a partir del archivo del
blog Dermapixel (\url{https://dermapixel.com}), mantenido desde
hace más de quince años por la Dra.~R.~Taberner, dermatóloga
clínica del Hospital Universitario Son Llàtzer. El archivo original
recoge casos clínicos comentados con fines docentes y representa la
materia prima a partir de la cual se ha estructurado el dataset
documentado en este anexo.

La construcción del dataset, el pipeline de extracción y
estructuración, la definición de la ontología jerárquica L1/L2/L3 y
el particionado han sido desarrollados por Antonio Contestí Coll en
el marco del presente Trabajo de Fin de Grado y del trabajo
paralelo de modernización tecnológica del archivo Dermapixel
descrito en el capítulo~\ref{cap:intro}. El uso del archivo
Dermapixel como materia prima del dataset se ha realizado bajo
acuerdo previo con la Dra.~R.~Taberner, formalizado por escrito en
el marco de la colaboración entre ambos sobre la modernización del
archivo y sobre el proyecto institucional de teledermatología del
IB-Salut. La revisión experta de las etiquetas diagnósticas, del
mapeo a la ontología jerárquica y de los campos de validación
asociados (\texttt{label\_source}, \texttt{diagnosis\_source},
\texttt{rosa\_verified}) ha sido aportada por la Dra.~Taberner.

En consecuencia, el dataset se publica con autoría compartida
Taberner \& Contestí Coll, en este orden, dado que el archivo
Dermapixel del que deriva el dataset es propiedad intelectual de la
Dra.~Taberner y constituye el material original sin el cual el
dataset no existiría. La cita recomendada
(sección~\ref{sec:dataset-cita}) refleja esta autoría conjunta en
el mismo orden.

\section{Licencia de uso}
\label{sec:dataset-licencia}

DermapixelAI 1.0 se publica bajo la licencia \textbf{Creative
Commons Reconocimiento--NoComercial--CompartirIgual 4.0
Internacional (CC BY-NC-SA 4.0)}. Texto canónico de la licencia:
\url{https://creativecommons.org/licenses/by-nc-sa/4.0/deed.es}.

Esta licencia implica, en lo que respecta al uso del dataset:

\begin{itemize}
\item \textbf{Reconocimiento (BY).} Cualquier uso del dataset o de
obras derivadas exige la cita explícita del dataset original en la
forma indicada en la sección~\ref{sec:dataset-cita}.
\item \textbf{No comercial (NC).} El dataset no puede emplearse con
fines primordialmente comerciales. La investigación académica, el
uso docente y el desarrollo de prototipos en fase no comercial
quedan dentro del alcance de la licencia.
\item \textbf{Compartir Igual (SA).} Cualquier obra derivada del
dataset ---incluidas, por ejemplo, versiones ampliadas o subconjuntos
re-anotados--- deberá distribuirse bajo la misma licencia o una
licencia compatible que preserve estas condiciones.
\end{itemize}

La elección de CC BY-NC-SA 4.0 responde a un equilibrio entre la
voluntad de los autores de facilitar la reutilización académica
del material y la naturaleza clínica del archivo original, derivado
de casos asistenciales reales publicados con finalidad docente.

\section{Cita recomendada}
\label{sec:dataset-cita}

La cita en prosa recomendada para cualquier publicación que utilice
el dataset es:

\begin{quote}
Taberner, R. y Contestí Coll, A. (2026). \emph{DermapixelAI 1.0:
dataset dermatológico clínico en castellano construido sobre el
archivo Dermapixel}. Versión 1.0. Licencia CC BY-NC-SA 4.0.
\end{quote}

Entrada BibTeX equivalente (pendiente de completar con el DOI tras
la publicación efectiva):

\begin{verbatim}
@dataset{dermapixelai2026,
  author    = {Taberner, Rosa and Contestí Coll, Antonio},
  title     = {DermapixelAI 1.0: dataset dermatológico
               clínico en castellano construido sobre
               el archivo Dermapixel},
  year      = {2026},
  version   = {1.0},
  license   = {CC BY-NC-SA 4.0},
  doi       = {pendiente},
  url       = {pendiente},
}
\end{verbatim}

\section{Disponibilidad y distribución}
\label{sec:dataset-disponibilidad}

La publicación efectiva del dataset en un repositorio externo está
pendiente de decisión en la fecha de cierre del presente trabajo.
La plataforma de distribución, la asignación del identificador
persistente (DOI) y la disponibilidad pública se documentarán en el
campo correspondiente de la cita BibTeX cuando se haya realizado el
depósito. Mientras tanto, las consultas sobre acceso al dataset
deben dirigirse al contacto indicado en la
sección~\ref{sec:dataset-mantenimiento}.

La estructura del paquete previsto se compone de:

\begin{itemize}
\item \texttt{dataset.csv}: fichero maestro con una fila por imagen
(1\,089 filas) y campos de modalidad, ruta, hash MD5, diagnóstico
L1/L2/L3, campos de validación experta, identificador de caso,
texto narrativo y partición del split.
\item \texttt{metadata/cases.csv}: una fila por caso clínico,
incluidos los casos sin imagen en el dataset
(\texttt{num\_images\_in\_dataset = 0}) por trazabilidad.
\item \texttt{metadata/ontology.csv}: vocabulario jerárquico
L1/L2/L3 con códigos SNOMED CT y CIE-10 asociados a cada entrada
L3.
\item \texttt{metadata/excluded\_images.csv}: registro de imágenes
descartadas durante el pipeline, con razón documentada.
\item \texttt{metadata/audit\_log.txt}: traza de auditoría de las
modificaciones aplicadas al dataset.
\item \texttt{splits/train.csv}, \texttt{splits/val.csv},
\texttt{splits/test.csv}: particiones canónicas case-aware.
\item \texttt{images/}: imágenes organizadas por modalidad
(\texttt{clinical/}, \texttt{dermoscopy/}, \texttt{histology/},
\texttt{ultrasound/}, \texttt{wood\_lamp/}).
\item \texttt{README.md}: documento que reproduce, en forma de
fichero plano, el contenido del presente anexo.
\item \texttt{LICENSE.txt}: texto canónico de la licencia
CC~BY-NC-SA~4.0 con cláusulas específicas de privacidad del
paciente y provenance clínico.
\item \texttt{CITATION.cff}: ficha de citación en formato Citation
File Format 1.2.0, con el campo \texttt{preferred-citation}
preparado para sustituir el DOI cuando Zenodo lo asigne.
\item \texttt{MANIFEST.sha256}: manifiesto global del paquete con
el hash SHA-256 de cada uno de los archivos distribuidos,
verificable con \texttt{sha256sum -c MANIFEST.sha256}.
\item \texttt{enriched\_captions/}: corpus enriquecido SpanDerm-CLIP
con captions sintéticas multi-perspectiva generadas por LLM
(sección~\ref{sec:spanderm-clip}) y \texttt{spanderm\_clip\_trilogy/}
con los resultados experimentales completos del experimento dual
encoder.
\end{itemize}

La integridad reproducible del paquete se verifica en dos niveles:
(i)~el manifiesto interno de cada partición o subconjunto
experimental, fijado en el \emph{checkpoint} correspondiente;
y~(ii)~el manifiesto global del paquete distribuido
(\texttt{MANIFEST.sha256}), que cubre los 1\,174 archivos del
\emph{bundle} y permite detectar cualquier alteración durante la
descarga o el almacenamiento intermedio. Ambos niveles son
independientes y se generan de forma determinista.

\section{Datasheet for Datasets}
\label{sec:dataset-datasheet}

Esta sección sigue el formato propuesto por~\cite{datasheets2018}
con las cuestiones más relevantes para el dataset.

\subsection*{Motivación}

\textbf{¿Para qué se creó el dataset?} El dataset se ha construido
para servir de material de investigación académica sobre
dermatología clínica en castellano, con foco en el aprendizaje
multimodal (imagen acompañada de texto clínico) y como base
estable para trabajos posteriores. No se ha creado como benchmark
de evaluación del presente TFG sino como contribución resultante
del proyecto a la comunidad investigadora.

\textbf{¿Quién lo financió?} La construcción del dataset se ha
realizado sin financiación específica, en el marco del presente
Trabajo de Fin de Grado, desarrollado en la Escola Politècnica
Superior de la Universitat de les Illes Balears, y de la colaboración
con la Dra.~R.~Taberner sobre la modernización del archivo
Dermapixel.

\subsection*{Composición}

\textbf{¿Qué representa cada instancia?} Cada instancia es una
imagen dermatológica asociada a un caso clínico publicado en el
archivo Dermapixel.

\textbf{¿Cuántas instancias hay?} 1\,089 imágenes vinculadas a 672
casos con imagen sobre 698 casos catalogados en total.

\textbf{¿Es el dataset un subconjunto de uno mayor?} El dataset se
construye a partir del archivo público del blog Dermapixel
(\url{https://dermapixel.com}), que tiene una composición más amplia
en términos de publicaciones y elementos no clínicos (banners,
ilustraciones, fotografías de eventos) que se filtran durante la
construcción del dataset (sección~\ref{sec:anexo-pipeline} del
anexo~\ref{anexo:pipeline}).

\textbf{¿Qué información acompaña a cada imagen?} Modalidad,
diagnóstico L1/L2/L3, identificador de caso, texto narrativo
asociado en castellano, campos de validación experta y
particionado en train/val/test.

\textbf{¿Hay etiquetas o ground-truth?} Sí. El diagnóstico L3 está
mapeado a la ontología jerárquica con revisión experta. El campo
\texttt{rosa\_verified} marca las imágenes adicionalmente validadas
visualmente por la Dra.~Taberner.

\textbf{¿Hay información que se haya eliminado?} Se han eliminado
las imágenes asociadas a la pista de solución de los casos del blog
(susceptibles de inducir \emph{leakage}) y las imágenes no
dermatológicas detectadas durante la auditoría de modalidad.

\subsection*{Proceso de recogida}

El detalle del pipeline figura en la
sección~\ref{sec:anexo-pipeline} del anexo~\ref{anexo:pipeline}.

\subsection*{Preprocesamiento, limpieza y etiquetado}

Hash MD5 por imagen para integridad; deduplicación exacta; mapeo
ontológico con revisión experta; particionado case-aware. Detalle en
el anexo~\ref{anexo:pipeline}.

\subsection*{Usos previstos}

\textbf{¿Para qué tipo de tareas se ha diseñado?} Investigación
académica sobre dermatología clínica en castellano, en particular
los escenarios que requieren material multimodal imagen-texto en
español, la armonización ontológica entre datasets dermatológicos
heterogéneos, y la docencia universitaria sobre análisis de
imagen médica.

\textbf{¿Para qué tareas \emph{no} debe usarse?} El dataset no
constituye un dispositivo médico ni puede emplearse como base de
un sistema de diagnóstico autónomo destinado a uso clínico
directo. Cualquier aplicación clínica derivada del dataset debe
seguir los procesos regulatorios correspondientes y no puede
sustituir al juicio del especialista.

\subsection*{Distribución}

CC BY-NC-SA 4.0. Plataforma pendiente de decisión en la fecha de
cierre del trabajo.

\subsection*{Mantenimiento}

Versionado semántico. Política de errata reportable mediante el
contacto de la sección~\ref{sec:dataset-mantenimiento}.

\section{Uso responsable y limitaciones clínicas}
\label{sec:dataset-uso-responsable}

El dataset DermapixelAI 1.0 se distribuye con las advertencias
siguientes:

\begin{itemize}
\item \textbf{No es un dispositivo médico.} El dataset es material
de investigación académica. Cualquier prototipo desarrollado a
partir del dataset que aspire a uso clínico real debe seguir los
procedimientos regulatorios aplicables (marcado CE, evaluación
clínica, etc.).

\item \textbf{No para uso diagnóstico autónomo.} Las predicciones
generadas por modelos entrenados o evaluados sobre el dataset no
pueden sustituir al juicio del dermatólogo. El dataset se ha
diseñado para apoyar la investigación en sistemas de apoyo al
diagnóstico, no para reemplazar la consulta especializada.

\item \textbf{Limitaciones de cobertura.} El dataset presenta
limitaciones específicas (desbalance de modalidad, cola larga
diagnóstica, cobertura ontológica incompleta, cobertura parcial
por subcategoría en los splits) documentadas en el
anexo~\ref{anexo:pipeline}. Cualquier publicación derivada debe
declarar estas limitaciones cuando reporte cifras de rendimiento.

\item \textbf{Generalización geográfica.} El dataset proviene del
archivo Dermapixel, alimentado por casos atendidos en el HUSLL en
Mallorca, España. La validez de los resultados sobre poblaciones
de fototipo, geografía o circuito asistencial distintos requiere
verificación independiente.
\end{itemize}

\section{Privacidad y consentimiento}
\label{sec:dataset-privacidad}

El archivo Dermapixel es un blog público dirigido a la formación
de profesionales sanitarios. Las imágenes publicadas en el blog
proceden de casos atendidos por la Dra.~R.~Taberner en su práctica
clínica y se publicaron originalmente con consentimiento implícito
en el marco del compromiso docente del blog y de la práctica
asistencial habitual sobre material clínico anonimizado para uso
formativo. El dataset DermapixelAI 1.0 hereda esta naturaleza de
material docente público.

El dataset no contiene metadatos identificativos directos de
pacientes (nombres, fechas de nacimiento, números de historia
clínica, ubicación geográfica concreta). Las imágenes están
recortadas a la lesión cuando procede; las fotografías macroscópicas
preservan únicamente la región anatómica relevante.

Cualquier reusuario del dataset se compromete, en el momento de
aceptar la licencia, a no intentar reidentificar a los pacientes
representados en las imágenes ni a usar el material con finalidad
distinta a la investigación académica o docente.

\section{Mantenimiento y contacto}
\label{sec:dataset-mantenimiento}

\textbf{Política de versionado.} El dataset se versiona siguiendo
una numeración semántica MAJOR.MINOR. Cambios en el conjunto de
imágenes incluidas o en la ontología canónica incrementan el
campo MAJOR. Correcciones de etiquetado, normalización de
campos o ampliaciones de los campos de validación experta
incrementan el campo MINOR.

\textbf{Reporte de errata.} Cualquier discrepancia detectada entre
las etiquetas asignadas y el diagnóstico correcto, o cualquier
problema de integridad, puede reportarse al contacto indicado a
continuación. Las erratas confirmadas se incorporan en la
siguiente versión MINOR del dataset.

\textbf{Contacto.}

\begin{itemize}
\item \textbf{Propietaria del archivo original y revisión experta
clínica:} Dra.~Rosa Taberner
\item \textbf{Autor técnico del dataset:} Antonio Contestí Coll
\item \textbf{Afiliación principal:} Hospital Universitario Son
Llàtzer y Universitat de les Illes Balears (Escola Politècnica
Superior)
\item \textbf{Canal de contacto preferente:} mediante el repositorio
oficial del dataset una vez publicado; mientras tanto, a través del
autor técnico del dataset, Antonio Contestí Coll.
\end{itemize}
